\documentclass[11pt]{article}
\usepackage[margin=2.5cm]{geometry}
\usepackage{amsmath,amssymb}
\usepackage{booktabs}
\usepackage{hyperref}
\title{Disproof of TLMC Conjecture 00000000570}
\author{TLMC Falsification Package}
\date{October 2, 2026}

\begin{document}
\maketitle

\begin{abstract}
Conjecture 00000000570 asserts that for every acyclic cluster algebra the
width of the $x$-degree support of the Laurent expansion in the initial
cluster (the Chebyshev-normalized invariant) is at most $2\cdot\mathrm{rank}$
(``locality of degrees'').  We disprove the conjecture with the Kronecker
cluster algebra (rank $2$, exchange matrix $B=\bigl[\begin{smallmatrix}0&2\\-2&0\end{smallmatrix}\bigr]$,
affine type $\widetilde A_1$, acyclic): the support width of the cluster
variable $x_6$ is $6 > 4 = 2\cdot\mathrm{rank}$, and the widths
$0,2,4,6,8,10,12,\dots$ of $x_3,x_4,x_5,\dots$ grow without bound.
All numbers are machine-verified: an independent pure-Python recomputation
(\texttt{reproduce.py}) and a core-Lean~4 formalization with zero axioms and
zero \texttt{sorry} (\texttt{lean4/}).
\end{abstract}

\section{The counterexample}

Let $x_1,x_2$ be the initial cluster of the Kronecker cluster algebra
(quiver $1 \rightrightarrows 2$, acyclic, rank $2$).  Alternating mutations
produce cluster variables satisfying the rank-2 exchange recurrence
\[
x_{n+1}\,x_{n-1} \;=\; x_n^2 + 1, \qquad x_3 = \frac{x_2^2+1}{x_1}.
\]
By the Laurent phenomenon each $x_n$ is a Laurent polynomial in $x_1,x_2$;
we compute these expansions exactly (greedy Laurent-polynomial division,
verified by remultiplication).  Writing the support of $x_n$ as a set of
exponent pairs $(a,b)\in\mathbb{Z}^2$, the $x_1$-exponent $a$ spans exactly
the interval $[-(n-2),\,n-4]$:

\begin{table}[h]
\centering
\begin{tabular}{rrll}
\toprule
$n$ & \#terms & $x_1$-range of support & width vs.\ $2\cdot\mathrm{rank}=4$ \\
\midrule
3 & 2  & $[-1,-1]$ & $0$ (ok) \\
4 & 4  & $[-2,0]$  & $2$ (ok) \\
5 & 7  & $[-3,1]$  & $4$ (ok) \\
6 & 11 & $[-4,2]$  & $\mathbf{6 > 4}$ \textbf{(violation)} \\
7 & 16 & $[-5,3]$  & $8 > 4$ \\
8 & 22 & $[-6,4]$  & $10 > 4$ \\
9 & 29 & $[-7,5]$  & $12 > 4$ \\
\bottomrule
\end{tabular}
\caption{Support widths along the Kronecker chain.  The values
$0,2,4,6,8,10,12$ for $x_3,\dots,x_9$ match the recorded attack exactly.}
\end{table}

Both endpoints of the strip are attained: for $x_6$,
\[
x_6 \;=\; x_1^{-4}x_2^{-3} + 4x_1^{-4}x_2^{-1} + 6x_1^{-4}x_2 + 4x_1^{-4}x_2^3
+ x_1^{-4}x_2^5 + 3x_1^{-2}x_2^{-3} + 6x_1^{-2}x_2^{-1} + 3x_1^{-2}x_2
+ 3x_2^{-3} + 2x_2^{-1} + x_1^{2}x_2^{-3},
\]
with the corner coefficients equal to $1$, so the width is exactly
$6 > 4$.  The growth is unbounded: $\mathrm{width}(x_n) = 2(n-3)$; the
accompanying script verifies this to $n=25$ (width $44$).  The inductive
mechanism is that the leading $x_1$-exponent of
$x_{n+1}=(x_n^2+1)/x_{n-1}$ is twice the top exponent of $x_n$ minus the
bottom exponent of $x_{n-1}$, widening the interval by one on each side per
step.  All other natural spreads ($x_2$-exponent, total degree
$x_1x_2$, and $x_1x_2^{-1}$-degree) also grow by $2$ per step, so the
disproof does not depend on the width convention.

\section{Why the numbers are the cluster variables}

The script asserts the seven Laurent-polynomial identities
$x_{n-1}x_{n+1} = x_n^2+1$ ($n=2,\dots,8$) exactly, and every division is
verified by remultiplication.  Since $\mathbb{Z}[x_1^{\pm1},x_2^{\pm1}]$ is
an integral domain, the Laurent polynomial $q$ with $q\cdot x_{n-1} = x_n^2+1$
is unique; hence the verified tables \emph{are} the Laurent expansions of
the cluster variables $x_1,\dots,x_9$.  The Kronecker quiver is acyclic and
stays acyclic under mutation, so the counterexample lies inside the
conjecture's hypothesis, and $\mathrm{width}(x_6)=6>2\cdot\mathrm{rank}=4$
falsifies it.

\section{Machine verification}

\textbf{Lean 4} (\texttt{lean4/}, core Lean only, no Mathlib): the tables
\texttt{TBL1}--\texttt{TBL9} are verified against all seven exchange relations
coefficientwise (\texttt{relz2}--\texttt{relz8}), the strip bounds and attained
endpoints give the exact widths
$0,2,4,6,8,10,12$ (\texttt{width3}--\texttt{width9}), and
\texttt{disproof\_00000000570} packages the violations
$\mathrm{width}(x_6)=6>4$ and $\mathrm{width}(x_9)=12>4$.
\texttt{Check.lean} audits \texttt{\#print axioms} for every declaration:
all report ``does not depend on any axioms'' (zero axioms, zero
\texttt{sorry}; \texttt{decide} only on closed propositions).

\textbf{Python} (\texttt{reproduce.py}, no dependencies): recomputes
$x_3,\dots,x_{25}$ from scratch, verifies all divisions and the seven
exchange identities, prints the width table and the $b=1$ boundary
comparison, and emits the Lean tables (byte-identical to
\texttt{lean4/Main.lean}).

\section{Boundary}

\begin{itemize}
\item \emph{Rank $1$}: the width is $0 \le 2\cdot\mathrm{rank}$; holds.
\item \emph{Rank $2$, finite type $A_2$} ($b=1$): widths stay $\le 1 \le 4$;
      holds.
\item \emph{Rank $2$, affine $\widetilde A_1$} ($b=2$, this attack): fails
      from $x_6$ on, unboundedly.
\item \emph{Rank $R \ge 3$ acyclic}: the product Kronecker $\times
      A_1^{R-2}$ is acyclic of rank $R$ with the same support widths, so the
      bound fails for every $R \ge 2$.
\item \emph{Chebyshev normalization}: normalizing cluster variables by
      Laurent monomials translates supports uniformly and leaves every
      width invariant, so the attack survives the normalization.
\end{itemize}

\end{document}
